% Lösungen Einheit 4 von 4 – Ausklammern (zusammenbau v0.1)
\einheitenkopf{Einheit 4 von 4 · Ausklammern}

%% TODO Lösung der Erklärzeile 20f) fehlt in der Bank
\erg{20}{a) $4 \cdot (x + 2)$ \quad b) $5 \cdot (a + 3)$ \quad c) $7 \cdot (b + 2)$ \quad d) $3 \cdot (z + 7)$ \quad e) $3 \cdot (3y + 4)$ \quad f) \ldots \quad g) $x \cdot (x + 5)$ \quad h) $3x \cdot (2x + 5)$ \quad i) $6 \cdot (x + 1)$ \quad j) $3 \cdot (x^2 + 2x + 5)$ \quad k) $6 \cdot (x + 5)$; Probe: $6 \cdot x + 6 \cdot 5 = 6x + 30$ \quad l) Die 20 wurde nicht durch 5 geteilt, die Probe ergibt $5x + 100$. Richtig: $5x + 20 = 5 \cdot (x + 4)$. \quad m) $3a \cdot (a + 2b + 3)$; Probe: $3a \cdot a + 3a \cdot 2b + 3a \cdot 3 = 3a^2 + 6ab + 9a$}
\erg{21}{a) Das Minus ging verloren, die Probe ergibt $12x + 8$. Richtig: $12x - 8 = 4 \cdot (3x - 2)$.}
